\subsection{Edges and minimal vertex covers as pruning supports}
\label{app:core-graph-prunings}

Suppose the first list of minimal supports consists of the edges of a
finite simple graph $G$ with at least one edge. Its minimal transversals
are the minimal vertex covers. Call $G$ realizable when one word has
every $+$-pruning supported on a union of edges and every $-$-pruning
supported on a union of minimal vertex covers. These are the two
union-closed families of the preceding subsection. In a bipartite graph
realizability is equivalent to distributivity of the lattice of unions
of neighborhoods. A different criterion applies to all graphs without
induced four-cycles.

This asks which words compute the Boolean function
$f_G=\bigvee_{uv\in E(G)}(x_u\land x_v)$ using only essential
inputs after any restriction fixing inputs to zero, or any restriction
fixing inputs to one. Here $\land=\star_+$ and $\lor=\star_-$.
Indeed, the Boolean minterms of a word are its minimal $+$-pruning
supports, and its minimal zero certificates are its minimal $-$-pruning
supports. The two lists determine one another by transversality, so a
word meeting the two prescribed bounds has exactly the displayed Boolean
function: its minterms contain edges and its zero certificates contain
vertex covers. Fix all inputs outside $U$ to zero and simplify the
constants in the word. The surviving inputs are the union of all
$+$-pruning supports contained in $U$, by induction on the gates.
The essential inputs of the restricted Boolean function are the union
of its minterms contained in $U$. These sets agree for every $U$
precisely when each pruning support is a union of minterms, as one
sees by taking $U$ to be that support. Fixing inputs to one gives the
dual assertion. Restrictions fixing both truth values at once give
the cograph criterion of \cref{core:thm:cograph-extension}.

Write $N(x)$ for the open neighborhood and $N[x]=N(x)\cup\{x\}$.
A nonempty set $H$ is a union of edges precisely when each of its
vertices has a neighbor in $H$. It is a union of minimal covers
precisely when it is a vertex cover and each $x\in H$ has a neighbor
$w$ with $N(w)\subseteq H$. Indeed, if a minimal cover $C\subseteq H$
contains $x$, some neighbor $w$ is outside $C$, so $N(w)\subseteq C$.
Conversely extend $(V\setminus H)\cup\{w\}$ to a maximal independent
set. Its complement is a minimal cover inside $H$ containing $x$.

\begin{lemma}[An obstruction at an edge]\label{core:lem:graph-edge-cut}
Let $uv$ be an edge. For $s=u,v$, let $T_s$ be an independent set
outside $N[s]$ such that each neighbor of $s$ has a neighbor in $T_s$.
If $(T_u\cap N(v))\cup(T_v\cap N(u))$ is independent, then $G$
is unrealizable.
\end{lemma}
\begin{proof}
Let $P$ be the complement of the displayed independent set and
$Q=\{u,v\}$. Every edge meets $P$, and every vertex cover meets $Q$,
so any realizing word passes the tests $+:P$ and $-:Q$. Keep the two
$-$-omissions $s\Rightarrow\text{hit }T_s$ for $s=u,v$.
They hold in every union of minimal covers: if a minimal cover
contains $s$, some neighbor of $s$ is outside the cover and has
all its neighbors in it, including one in $T_s$. Keep also the
$+$-omissions $j\Rightarrow\text{hit }N(j)$ for
$j\in (T_u\cup T_v)\cap P$.

Consider a shortest word meeting these omissions and tests. At a
$+$-root both children retain $-:Q$ and the $-$-omissions. Each
selected $N(j)$ misses $Q$, so neither child can pass $+:N(j)$.
Both must therefore obey all the $+$-omissions, and one retains $+:P$.
At a $-$-root both children retain $+:P$ and the $+$-omissions.
If either passed $-:T_s$ it would have a $+$-pruning inside $T_s$
that meets $P$. Its vertex $j\in P$ would need a neighbor in that
pruning, contrary to independence. Both children therefore obey the
$-$-omissions, and one retains $-:Q$. A leaf meeting both tests must
be $u$ or $v$ and fails its own $-$-omission. This excludes every word.
\end{proof}

A vertex $x$ is \emph{shedding} when every independent set outside
$N[x]$ can be extended by a neighbor of $x$
\cite[Definition~5]{BiyikogluCivan2014}. Thus $x$ fails to be shedding
exactly when a set $T_x$ as in \cref{core:lem:graph-edge-cut} exists.

\begin{theorem}[Graphs without induced four-cycles]
\label{core:thm:shedding-prunings}
If $G$ has no induced four-cycle, it is realizable exactly when its
shedding vertices cover its edges. The sufficient direction holds for
every graph.
\end{theorem}
\begin{proof}
If $uv$ has two non-shedding endpoints, choose $T_u,T_v$ as above.
An edge between $T_u\cap N(v)$ and $T_v\cap N(u)$ would complete $uv$
to an induced four-cycle. The lemma therefore applies.

For sufficiency, write $\lor=\star_-$ and $\land=\star_+$ so that a
$+$-pruning selects a disjunct and keeps both conjuncts. Use the word
\[
  \bigvee_{v\text{ shedding}}
     \left(x_v\land\bigvee_{w\in N(v)}x_w\right)
\]
over the nonisolated shedding vertices. Every $+$-pruning yields an
edge. A $-$-pruning chooses either $\{v\}$ or all of $N(v)$ at each
center; their union $H$ covers all edges. For $x$ from a chosen $N(v)$,
the neighbor $v$ has $N(v)\subseteq H$. For a chosen center $x$,
either it has a neighbor in $V\setminus H$, whose entire neighborhood
lies in $H$, or $V\setminus H$ is independent outside $N[x]$.
In the latter case shedding gives a neighbor whose neighborhood
misses $V\setminus H$. Thus $H$ is a union of minimal covers.
\end{proof}

In a graph with no cycle of length at most five, a nonisolated vertex
is shedding exactly when it is adjacent to a leaf. Otherwise for each
neighbor choose a vertex at distance two through it. The chosen set
is independent, since an edge between two choices from different
neighbors would form a five-cycle. It dominates the neighbors and
lies outside the closed neighborhood. Checking coverage by the
neighbors of leaves takes linear time in the size of the graph.
In particular the five-vertex path is realizable and the six-vertex
path is unrealizable.

\begin{corollary}[Complexity without triangles or squares]
\label{core:cor:shedding-pruning-complexity}
Recognizing realizability from the graph is coNP-complete on graphs
with no triangles or four-cycles.
\end{corollary}
\begin{proof}
By \cref{core:thm:shedding-prunings}, a certificate of failure is
an edge with two non-shedding endpoints, together with their two
independent sets of witnesses.

For hardness take a CNF formula with a nonempty list of nonempty
clauses and distinct variables in each clause
\cite{Cook1971}. We can assume every variable occurs at most three
times: split its $m$ occurrences into fresh variables
$x_1,\ldots,x_m$ and add the implications $x_i\to x_{i+1}$ in a
cycle when $m>1$. This preserves satisfiability and the clause
condition. Each new variable has at most one occurrence with its
less frequent sign.

Make one vertex $o$ adjacent to all clause vertices. Join each
clause vertex to its own occurrence vertices, and join occurrence
vertices when they give opposite signs to the same variable.
Attach a leaf to every occurrence vertex, and add a path $o-d-p-q$
with $q$ of degree one. The conflict graph on the occurrences is a
disjoint union of stars. A triangle through a clause would require
opposite signs of the same variable in that clause. A four-cycle
through $o$ would require an occurrence shared by two clauses.
A four-cycle avoiding $o$ and meeting a clause would require either
a shared occurrence or two occurrences of one clause with a
common conflict neighbor. Both are impossible, since occurrences
belong to single clauses and each variable appears at most once
per clause. Leaf attachments and the added path create no cycles.

An independent set outside $N[o]$ dominates $N(o)$ exactly when it
contains $p$ and contains, for every clause, some occurrence in
that clause. Such choices are consistent precisely when the
formula is satisfiable. Thus $o$ is shedding exactly in the
unsatisfiable case. The vertex $d$ is non-shedding: $q$ and any
clause vertex form an independent set outside $N[d]$ dominating
$N(d)$. Every edge outside $o$ has an endpoint adjacent to a leaf,
either an occurrence vertex or $p$. Therefore the shedding
vertices cover all edges when $o$ is shedding, and miss $od$ when
it is non-shedding.
\end{proof}

\begin{theorem}[Bipartite graphs]\label{core:thm:bipartite-prunings}
Let $G$ have parts $X,Y$. Call $ba$, with $b\in X$ and $a\in Y$,
\emph{bisimplicial} when every vertex of $N(a)$ is adjacent to every
vertex of $N(b)$. The following conditions are equivalent:
\begin{enumerate}
\item $G$ is realizable.
\item The rectangles $N(a)\times N(b)$ of its bisimplicial edges cover $E(G)$.
\item The lattice $L=\{\bigcup_{x\in S}N(x):S\subseteq X\}$, ordered
by inclusion, is distributive.
\end{enumerate}
\end{theorem}
\begin{proof}
For a given edge $xy$ take an edge $uv$, with $u\in X$, satisfying
$N(u)\subseteq N(x)$, $N(v)\subseteq N(y)$ and minimizing
$|N(u)|+|N(v)|$. Then $u$ is adjacent to $y$ and $v$ to $x$.
If $uv$ is bisimplicial, its rectangle covers $xy$. Otherwise pick
nonadjacent $d\in N(v)$, $e\in N(u)$. For each $w\in N(u)$
missing $d$, minimality gives some $d_w\in N(w)\setminus N(v)$:
otherwise $uw$ improves on $uv$, since $d\in N(v)\setminus N(w)$.
The set $T_u$ of $d$ and these $d_w$ is independent on the $X$ side,
avoids $u$, and dominates $N(u)$, with $T_u\cap N(v)=\{d\}$.
Construct $T_v$ symmetrically with $T_v\cap N(u)=\{e\}$.
\Cref{core:lem:graph-edge-cut} rules out realizability. This proves
necessity of the rectangle cover.

Given the cover, take the disjunction over its rectangles of
$(\bigvee_{z\in N(a)}x_z)\land(\bigvee_{w\in N(b)}x_w)$.
A $+$-pruning selects an edge. A $-$-pruning chooses one entire
shore from every rectangle. Their union $H$ covers the graph, and
any selected vertex belongs to some $N(s)\subseteq H$, giving its
required neighbor $s$. Thus the word realizes $G$.

Write $p(S)=\bigcup_{x\in S}N(x)$ and let $s(U)$ contain those
$b\in X$ for which $ba$ is bisimplicial for some $a\in U$.
Both maps preserve unions and the rectangle cover gives $p(s(U))=U$
for $U\in L$: if $a\in U=p(S)$ then every $N(b)$ selected by $a$
is contained in $p(S)$; for each edge from $S$ the rectangle covering
it selects a row containing its other endpoint. Now $U\subseteq V\cup W$
in $L$ gives
\[
 U=p(s(U)\cap s(V))\ \cup\ p(s(U)\cap s(W)).
\]
The two terms lie below $U\wedge V$ and $U\wedge W$ respectively.
Applying this with $U=Z\wedge(V\cup W)$ gives distributivity.

Conversely, in a finite distributive lattice each nonzero join-irreducible
$J$ is join-prime, and each element is a join of these. Such a $J$ is
a row $N(b)$ since the rows generate $L$. The union of all rows not
containing $J$ misses some $a\in J$ by join-primality. Thus $a\in N(x)$
exactly when $J\subseteq N(x)$. The edge $ba$ is bisimplicial. For
any $y\in N(x)$ take an irreducible $J\subseteq N(x)$ containing $y$;
its rectangle covers $xy$.
\end{proof}

The rectangle criterion applies more generally to graphs containing
neither a triangle nor a five-cycle. Use oriented edges $xy,uv$ in
the same minimization. Any two vertices at distance two from $u$
are nonadjacent: an edge between them, together with the two
length-two paths from $u$, would give a closed walk of length five
and hence a triangle or a five-cycle. The members chosen for $T_u$
are at distance two from $u$ (triangles are excluded), so $T_u$ is
again independent. The rectangle construction uses only the complete
adjacency between its shores and works in any graph.

Equivalently the bipartite incidence relation $E$ admits a relation
$D\subseteq Y\times X$ with $EDE=E$, using relational composition.
For a rectangle cover take the reversed bisimplicial edges as $D$.
Conversely such a $D$ makes $U\mapsto D[U]$ a union-preserving section
of $p$ on $L$, so the same distributivity argument applies.
The relation between regularity and distributivity is classical
\cite{Zaretskii1962}. Here it determines exactly when a common word
exists. For the six-cycle the neighborhood unions are the five-element
diamond, and for the six-vertex path they are the five-element pentagon.

For a general graph, a set $S$ is a \emph{robust shore} if, in every
vertex cover $H$ containing $S$, each vertex of $S$ belongs to a minimal
cover inside $H$. Equivalently, for each $s\in S$, the set $S$ meets
every independent set outside $N[s]$ that dominates $N(s)$.
Indeed $V\setminus H$ is independent. If $s$ has a neighbor there,
that neighbor's entire neighborhood lies in $H$. Otherwise the
criterion for a minimal cover through $s$ fails exactly when every
neighbor of $s$ has a neighbor in $V\setminus H$. Every full
neighborhood is robust. The rectangle construction above works for
any edge cover by complete rectangles with robust shores.

\begin{theorem}[Realization with edge supports]
\label{core:thm:edge-support-realization}
A graph has a realizing word whose $+$-prunings are all individual
edges exactly when it has a cover by complete rectangles with
robust shores.
\end{theorem}
\begin{proof}
The rectangle construction gives the required word. For the converse
fix an edge $uv$ and take a shortest word with all $+$-supports edges,
one equal to $\{u,v\}$, and all $-$-supports robust. A realizing word
satisfies these requirements: each edge is one of its minimal
$+$-supports. In this local problem a $-$-support may fail to
cover the graph.

At a disjunction, the child supplying $\{u,v\}$ still has all
$+$-supports edges. Its $-$-supports are robust. Indeed, if one
contained a vertex $j$ and missed an independent set $I$ outside
$N[j]$ dominating $N(j)$, the sibling, with all $+$-supports edges,
would be false on $I$. It would have a $-$-support disjoint from
$I$, so their union would disprove robustness at the root.
Minimality excludes this disjunction.

At a conjunction both children have robust $-$-supports. Each
child's $+$-supports are singletons or edges, since joining one
with any support of the sibling yields an edge. Each child has
a singleton support, as otherwise it would satisfy all the local
requirements by itself. Let $L,R$ collect these singleton labels
in the two children. The sets $L,R$ are robust. For instance an
independent set $I$ missing $L$ contains no $+$-support of the
left child. A $-$-support disjoint from $I$ must contain every
label of $L$, so $I$ cannot witness a failure of robustness for
any such label. Every pair in $L\times R$ is an edge.

If both shores meet $\{u,v\}$ their rectangle covers $uv$.
Otherwise the child whose shore misses $\{u,v\}$ supplies the
whole edge as a $+$-support. Then the other shore is a nonempty
robust subset of $\{u,v\}$. Each of its vertices is shedding,
since a witness of non-shedding outside its closed neighborhood
would miss both endpoints. Its star supplies a robust rectangle
covering $uv$. Repeat for each edge.
\end{proof}

\begin{proposition}[A graph requiring more than robust rectangles]
\label{core:prop:nonrectangular-prunings}
The graph on $u,v,a,b,x,y,z$ with edges
\[
 uv,\ ua,\ ub,\ vx,\ vy,\ ax,\ by,\ xy,\ za,\ zb
\]
is realizable, yet the edge $uv$ belongs to no complete rectangle
with robust shores. Hence every realizing word has a $+$-pruning
with at least three labels.
\end{proposition}
\begin{proof}
On a shore containing $u$ in such a rectangle, every vertex is
in $N(v)=\{u,x,y\}$. The independent sets $\{z,x\}$ and $\{z,y\}$
are outside $N[u]$ and dominate $N(u)$, so the shore must contain
both $x$ and $y$. The opposite shore contains $v$ and must meet
$\{a,b\}$, an independent set outside $N[v]$ dominating $N(v)$.
But $a$ misses $y$ and $b$ misses $x$.

Write the vertex names also for the corresponding leaves. Use
\[
 (z\lor u)\land(a\lor b)\quad\lor\quad
 \bigl(y\lor((x\lor u)\land(v\lor a))\bigr)
 \land
 \bigl(x\lor((y\lor u)\land(v\lor b))\bigr),
\]
with the two conjuncts on the right grouped as one disjunct.
The first disjunct supplies the four edges between $\{z,u\}$ and
$\{a,b\}$. The latter shore is $N(z)$. A cover containing $\{z,u\}$
must contain $x$ or $y$, and hence contain $N(a)$ or $N(b)$,
which witnesses robustness of $\{z,u\}$.

In the second disjunct a $+$-pruning joins two supports. The first
is $\{y\}$ or an edge from $\{x,u\}\times\{v,a\}$; the second
is $\{x\}$ or an edge from $\{y,u\}\times\{v,b\}$.
The two singletons give $xy$, and either singleton has a neighbor
in every opposing edge. Using these choices gives each of
$uv,vx,vy,ax,by,xy$ as a support. The $-$-prunings are exactly
$N(v),N(x),N(y)$. Thus every $+$-pruning of the whole word is a
union of edges, and every edge occurs as such a support. Every
$-$-pruning therefore covers the graph and is a union of robust
shores, each of whose vertices lies in a minimal cover inside it.
\end{proof}
